%% ============================================================
%%  Lecture 1 -- Monday, 6 July 2026
%%  Subfile of main.tex: compiles standalone OR as part of the book.
%%  Built from sources/2026-07-06/ (slide-transcript.md, outline.md).
%% ============================================================
\documentclass[main.tex]{subfiles}

\begin{document}

\beginlecture{1}{Monday, 6 July 2026}%
  {From Data to Decisions: Statistical Models, Loss, and Risk}

\epigraph{Far better an approximate answer to the right question, which is often vague, than an
exact answer to the wrong question, which can always be made precise.}%
{John W. Tukey, ``The Future of Data Analysis'' (1962)}

\begin{lecturecontents}{Contents of this lecture}
  \item \hyperlink{1.1}{1.1\quad Probability Review I: Randomness, Random Variables, Distributions, and Realizations}
  \item \hyperlink{1.2}{1.2\quad The Applied-Statistics Pipeline}
  \item \hyperlink{1.3}{1.3\quad The Statistical Model and the Analytic Goal}
  \item \hyperlink{1.4}{1.4\quad Probability Review II: Expectation, Conditional Expectation, and Iterated Expectations}
  \item \hyperlink{1.5}{1.5\quad Decision Theory: Actions, Loss, and Risk}
  \item \hyperlink{1.6}{1.6\quad Bayesian Decision Theory}
  \item \hyperlink{1.7}{1.7\quad The Cheat Sheet: Decision Theory in Eleven Items}
  \item \hyperlink{1.8}{1.8\quad Estimation, Testing, and the Clinical Trial}
\end{lecturecontents}

\bigskip
\noindent\textbf{Overview.} The course opens where applied statistics leaves off: a study has a
goal, data, and software output---and none of those, by themselves, say anything about reality.
This lecture builds the bridge. After a review of random variables, distributions, and
realizations (and, later, of expectation up through iterated expectations), the data's unknown
joint distribution is organized into a \emph{statistical model}, the study's purpose becomes an
\emph{analytic goal} expressed in the model's parameters, and the choice of method becomes a
decision problem: an action space, a loss function, and---as the figure of merit---the
\emph{risk function}, the expected loss as a function of the parameter. Admissibility and the
Bayes principle give the two classical answers to what ``good'' means, and both are put to work
on a two-arm clinical trial in which two estimators of a cure-rate difference compete. The
companion text is Hoel--Port--Stone, Chapter~1 (\HPS{1}).

\clearpage
% ============================================================
\sub{1.1}{Probability Review I: Randomness, Random Variables, Distributions, and Realizations}

\begin{ObjL}{Concept}{1.1.1}{The point of statistical theory}
``The point of statistical theory is to inform the application of statistics.'' The lecture
therefore begins with applied statistics---what a data analysis is \emph{for}---and works
backward to the theory that disciplines it.
\end{ObjL}

\begin{ObjL}{Remark}{1.1.2}{Randomness---does it exist?}
The review opens with a philosophical aside: whether randomness ``exists'' in the world or only
describes our ignorance is left unresolved---and is beside the point. Probability is the
\emph{model} we adopt for quantities whose values we cannot predict; everything in this course
is a statement within that model.
\end{ObjL}

\begin{Obj}{Definition}{1.1.3}{Sample space}
The \emph{sample space} $S$ is the set of all possible outcomes of the random phenomenon being
modeled. Events are subsets of $S$, and probabilities are assigned to events in accordance with
the axioms of probability.
\end{Obj}

\begin{Obj}{Definition}{1.1.4}{Random variable}
In the real world, a \emph{random variable} is a numerical quantity whose value is determined
by the outcome of a random phenomenon---a measurement not yet made. Mathematically, a random
variable is a (measurable) function
\[ X : S \to \mathbb{R}, \]
assigning a number $X(s)$ to each outcome $s\in S$. The data of a study are random variables;
so is anything computed from them.
\end{Obj}

\begin{Obj}{Definition}{1.1.5}{Distribution}
The \emph{distribution} of a random variable $X$ is the assignment
\[ A \;\longmapsto\; P(X\in A) \]
of a probability to each (measurable) set $A\subseteq\mathbb{R}$: the complete description, in
advance of observation, of how likely $X$ is to fall anywhere. It is determined by the
cumulative distribution function $F(x)=P(X\le x)$, and, in the discrete and continuous cases,
by a probability mass function or a density.
\end{Obj}

\begin{Obj}{Definition}{1.1.6}{Joint distribution, joint pmf, joint density}
The \emph{joint distribution} of random variables $X_1,\dots,X_n$ assigns a probability
$P\big((X_1,\dots,X_n)\in B\big)$ to each (measurable) $B\subseteq\mathbb{R}^n$; it carries
every dependence among the variables, and it determines---but is not in general determined
by---the marginal distribution of each $X_i$. In the discrete case it is given by the
\emph{joint probability mass function}
$p(x_1,\dots,x_n)=P(X_1=x_1,\dots,X_n=x_n)$; in the continuous case by a \emph{joint
density}\note{What a density \emph{is}: not a probability but a probability
\emph{rate}---probability per unit length (per unit area or volume, for joint densities).
Probabilities come only from integrating it over regions, so for small $dx$,
$f(x)\,dx \approx P(x\le X\le x+dx)$. Two consequences that trip people up: a density can
exceed $1$ (the uniform distribution on $[0,\tfrac12]$ has $f\equiv 2$ there), and every
single point carries probability zero ($\int_x^x f=0$), so for a continuous random variable
$f(x)$ is emphatically \emph{not} $P(X=x)$. The pmf, by contrast, \emph{is} a probability at
each point; the twin formulas---sums against $p$, integrals against $f$---conceal that
difference in kind.}
$f\ge0$ with
\[ P\big((X_1,\dots,X_n)\in B\big)\;=\;\int_B f(x_1,\dots,x_n)\,dx_1\cdots dx_n . \]
\end{Obj}

\begin{Obj}{Definition}{1.1.7}{Realization}
If the outcome of the experiment is $s\in S$, the \emph{realization} of the random variable $X$
is the number $x=X(s)$: the value actually observed. Random variables are written with capital
letters and their realizations in lower case. Your data set is a realization of the random
variables the model describes.
\end{Obj}

\begin{ObjL}{Example}{1.1.8}{The same coin flipped twice}
Flip the same coin twice, and let $X_1$ and $X_2$ be the head-indicators of the two flips. The
deck poses three questions.
\emph{Are the random variables the same?} No: $X_1$ and $X_2$ are different functions on the
sample space $S=\{HH,HT,TH,TT\}$---the first reads off the first coordinate, the second the
second.
\emph{Are the distributions the same?} Yes: each is Bernoulli with the same head-probability,
because it is the same coin.
\emph{Are the realizations the same?} Not necessarily: the observed pair may be $(1,0)$.
Distinct random variables can share one distribution and still realize differently---the three
notions of \xref{1.1.4}, \xref{1.1.5}, and \xref{1.1.7} are genuinely different.
\end{ObjL}

% ============================================================
\sub{1.2}{The Applied-Statistics Pipeline}

\begin{ObjL}{Concept}{1.2.1}{The goal of the study versus the data}
``If you are doing a statistical analysis, there should be something you are trying to
learn\dots\ NOT something about your data, something about what your data says about
reality.'' The goal of the study and the data are separate objects; no arrow connects them
directly.
\end{ObjL}

\begin{ObjL}{Concept}{1.2.2}{Methods and output}
Statistical methods are applied to the data---through a package (R, SAS, SPSS, Stata, \dots),
through modules (SciPy, NumPy, Pandas, \dots), ``or \dots\ AI?''---and ``the application of
methods produces output\dots''
\end{ObjL}

\begin{ObjL}{Concept}{1.2.3}{Interpretation}
``These results do not directly speak to the goal!'' They must be \emph{interpreted}, and it is
the interpretation of the results \emph{in light of where the data came from}---causality and
representativeness---that allows conclusions about reality.
\end{ObjL}

\begin{Fig}{1.2.4}{The applied-statistics pipeline}{What you know and what you don't know about
the phenomenon and the data collection feed the statistical model; the model and the goal of
the study produce the analytic goal; model and analytic goal select the method; the method
applied to the data produces results, which---interpreted in light of causality and
representativeness---yield conclusions. The dashed arrow is the feedback from results back into
the model. Recreated from the lecture deck (\xref{1.2.1}--\xref{1.2.3}; cf.\ \xref{1.3.1},
\xref{1.3.5}).}
\begin{tikzpicture}[x=1cm,y=1cm,
  every node/.style={font=\small},
  box/.style={draw, rounded corners=2pt, align=center, inner sep=4pt, minimum height=0.72cm,
    execute at begin node={\hyphenpenalty=10000\exhyphenpenalty=10000}},
  arr/.style={->, thick}]
  % top row: inputs to the model
  \node[box, draw=cuBlueDk, fill=cuBlueLt!35, text width=3.9cm] (know) at (0.1,6.3)
    {\textbf{What you know}\\[-1pt]{\footnotesize about the phenomenon and data collection}};
  \node[box, draw=figBad, fill=figBad!8, text width=3.9cm] (dont) at (4.5,6.3)
    {\textbf{What you don't know}\\[-1pt]{\footnotesize about the phenomenon and data collection}};
  \node[box, draw=cuBlueDk, fill=white, text width=3.3cm] (goal) at (8.6,6.3)
    {\textbf{The Goal of the Study}};
  % model row
  \node[box, draw=figLim, fill=figLim!10, text width=3.0cm] (model) at (1.9,4.6)
    {\textbf{Statistical Model}};
  \node[box, draw=figTerm, fill=figTerm!8, text width=3.3cm] (agoal) at (8.6,4.6)
    {\textbf{The Analytic Goal}};
  % method row
  \node[box, draw=figTerm, fill=figTerm!8, text width=2.2cm] (method) at (4.4,3.0)
    {\textbf{Method}};
  \node[box, draw=cuBlueDk, fill=cuBlueLt!45, text width=1.8cm] (data) at (8.3,3.0)
    {\textbf{Data}};
  % results row
  \node[box, draw=black!65, fill=black!4, text width=2.2cm] (results) at (6.3,1.5)
    {\textbf{Results}};
  \node[box, draw=figBad!80, fill=figBad!6, text width=3.4cm] (causal) at (10.9,1.5)
    {\textbf{Causality and Representativeness}};
  % interpretation
  \node[box, draw=figBad, fill=figBad!10, text width=2.6cm] (interp) at (8.6,0.0)
    {\textbf{Interpretation}};
  % arrows
  \draw[arr, cuBlueDk] (know.south) -- ([xshift=-0.9cm]model.north);
  \draw[arr, figBad!85!black] (dont.south) -- ([xshift=0.9cm]model.north);
  \draw[arr, figLim] (model.east) -- (agoal.west|-model.east);
  \draw[arr, cuBlueDk] (goal.south) -- (agoal.north);
  \draw[arr, figLim] (model.south) |- ([yshift=3.4pt]method.west);
  \draw[arr, figTerm] (agoal.south) .. controls (7.2,3.9) .. ([yshift=3.4pt]method.east);
  \draw[arr, figTerm] (method.south) |- ([yshift=2.8pt]results.west);
  \draw[arr, cuBlueDk] (data.south) |- ([yshift=2.8pt]results.east);
  \draw[arr, black!65] (results.south) |- ([yshift=2.8pt]interp.west);
  \draw[arr, figBad!80] (causal.south) |- ([yshift=2.8pt]interp.east);
  % dashed feedback: results -> model, up the left margin
  \draw[->, thick, dashed, figLim!80!black]
    (results.west) -- ++(-6.6,0) |- ([yshift=-1.4pt]model.west);
\end{tikzpicture}
\end{Fig}

% ============================================================
\sub{1.3}{The Statistical Model and the Analytic Goal}

\begin{Obj}{Definition}{1.3.1}{The statistical model}
If your data are random, they have a joint distribution (\xref{1.1.6}). What can you learn from
your data? ``The data! AND THE DISTRIBUTION OF THE DATA?! And that is all!'' If the study
has a goal, there must be something you do not know about that distribution---so more than one
distribution must be possible. \emph{The statistical model is the set of all possible
distributions of the data.} \HPS{1}
\end{Obj}

\begin{Obj}{Definition}{1.3.2}{Parameterization}
It is common to \emph{index} the family of distributions: a \emph{parameterization} assigns to
each value $\theta$ of a \emph{parameter space} $\Theta$ a distribution in the model, so the
model may be written
\[ \{\, p^{\theta}\ \text{or}\ f^{\theta} \;:\; \theta\in\Theta \,\} \]
with $p^{\theta}$ a pmf or $f^{\theta}$ a density for the data.
\end{Obj}

\begin{Obj}{Definition}{1.3.3}{The analytic goal}
The \emph{analytic goal} is an expression of the purpose of the study in terms of the
parameters of the model. The scientific question, which lives in the real world, is translated
into a statement about $\theta$.
\end{Obj}

\begin{ObjL}{Example}{1.3.4}{A two-arm clinical trial (the questions)}
A clinical trial compares a treatment arm with a placebo arm, measuring cure or failure for
each patient; the goal is to determine whether the treatment is efficacious. What is the
statistical model? What is the analytic goal? The formal answers are assembled in
\xref{1.8.4}, once decision theory has supplied the vocabulary.
\end{ObjL}

\begin{ObjL}{Concept}{1.3.5}{Where does the statistical model come from?}
From \emph{what you know} about the phenomenon and the data collection (which distributions are
plausible) and \emph{what you don't know} (which is why the model contains more than one
distribution)---the two inputs at the top of Figure~\xref{1.2.4}. The dashed feedback arrow
records that results can send you back to revise the model.
\end{ObjL}

\begin{ObjL}{Remark}{1.3.6}{What this class is about}
Mostly: how do we get from \emph{a model and an analytic goal} to \emph{a good method}---and
first, what do we mean by ``good''? In the deck's own strongest emphasis, these are ``an
answer, of sorts, to the question: what is it all about, extracting information from data?''
\end{ObjL}

% ============================================================
\sub{1.4}{Probability Review II: Expectation, Conditional Expectation, and Iterated Expectations}

The roadmap for the rest of the lecture: ``First, all the theory. Then a list of the essential
points. Then an illustration (the clinical trial).'' But first, some more probability review.

\begin{Obj}{Definition}{1.4.1}{Expectation}
Expectations belong to \emph{distributions}: a random variable has an expectation through its
distribution. For discrete $Y$ with pmf $p$ and for continuous $Y$ with density $f$,
\[ \E\{Y\} \;=\; \sum_{y} y\,p(y) \qquad\text{or}\qquad \E\{Y\} \;=\; \int y\,f(y)\,dy , \]
provided the sum or integral converges absolutely\note{\emph{Absolute} convergence means the
sum or integral of $|y|$ is finite---$\sum_y |y|\,p(y)<\infty$ or $\int|y|\,f(y)\,dy<\infty$,
i.e.\ $\E\{|Y|\}<\infty$---not merely that the signed sum happens to settle. The assumption
earns its keep twice. First, a distribution's values come in no privileged order, and by
Riemann's rearrangement theorem a conditionally convergent series can be reordered to reach
\emph{any} value---so without absolute convergence, ``the'' expectation would depend on
bookkeeping. Take $Y$ with $P\big(Y=(-1)^{k+1}2^k/k\big)=2^{-k}$ for $k=1,2,\dots$: summed in
the natural order, $\sum_k(-1)^{k+1}/k=\ln 2$, yet $\E\{|Y|\}=\sum_k 1/k=\infty$, and
rearranging the terms changes the answer---so such a $Y$ is assigned \emph{no} expectation.
Second, absolute convergence is precisely the license for the regrouping and
exchange-of-order steps in the proofs of \xref{1.4.4} and \xref{1.4.14} (for integrals, via
Fubini's theorem).}: the probability-weighted average of the
values of $Y$.\note{Four conflated words, disentangled. \emph{Expectation} is the number
defined here---a property of the distribution, fixed before any data are seen. \emph{Mean} is
its synonym when speaking of a distribution (a ``population mean''), but a \emph{sample mean}
is something else: it is an \emph{average}---the arithmetic operation $\frac1n\sum_i y_i$ on a
finite list of numbers, defined with no probability in sight. Averaging realized data produces
a random quantity that \emph{estimates} the expectation; only when every possible value is
equally likely does the expectation reduce to the plain average of those values. Finally,
``\emph{balancing point}'' is the physical reading of the same number: place the probability
as mass along the line, and $\E\{Y\}$ is its center of mass---the unique $c$ with
$\E\{Y-c\}=0$. It need not be a possible value: a fair die balances at $3.5$.}
\end{Obj}

\begin{Fig}{1.4.2}{The expectation as a balancing point}{Left: the fair die puts mass
$\tfrac16$ at each of $1,\dots,6$; the beam balances at $\E\{Y\}=3.5$, a value the die never
shows. Right: for a skewed density, probability is area, and the long right tail carries
enough of it to drag the balance point past the peak. In both cases $\E\{Y\}$ is the center
of mass of the distribution---the unique $c$ with $\E\{Y-c\}=0$ (\xref{1.4.1}).}
\begin{tikzpicture}[x=0.85cm,y=1.55cm]
  % ---- left panel: the fair die's pmf as point masses on a beam ----
  \draw[cuBlueDk] (0.4,0) -- (6.6,0);
  \foreach \x in {1,...,6}{
    \draw[figTerm, thick] (\x,0) -- (\x,0.42);
    \fill[figTerm] (\x,0.42) circle (2.6pt);
    \node[cuBlueDk, below, font=\footnotesize] at (\x,-0.02) {$\x$};
  }
  \node[figTerm, font=\footnotesize, left] at (0.9,0.42) {$p(y)=\tfrac16$};
  \draw[figLim, densely dashed] (3.5,0) -- (3.5,0.58);
  \fill[figLim] (3.5,-0.03) -- (3.3,-0.30) -- (3.7,-0.30) -- cycle;
  \node[figLim!80!black, font=\footnotesize, below] at (3.5,-0.34) {$\E\{Y\}=3.5$};
  % ---- right panel: a skewed density balancing past its peak ----
  \begin{scope}[shift={(8.6,0)}]
    \fill[figTerm!10] plot [smooth, tension=0.8] coordinates
      {(0.2,0) (0.7,0.30) (1.2,0.62) (1.7,0.72) (2.2,0.60) (3.2,0.36) (4.2,0.18) (5.2,0.08) (6.2,0.02)}
      -- (6.2,0) -- (0.2,0) -- cycle;
    \draw[figTerm, thick] plot [smooth, tension=0.8] coordinates
      {(0.2,0) (0.7,0.30) (1.2,0.62) (1.7,0.72) (2.2,0.60) (3.2,0.36) (4.2,0.18) (5.2,0.08) (6.2,0.02)};
    \draw[cuBlueDk] (0,0) -- (6.6,0);
    \draw[figLim, densely dashed] (2.55,0) -- (2.55,0.52);
    \fill[figLim] (2.55,-0.03) -- (2.35,-0.30) -- (2.75,-0.30) -- cycle;
    \node[figLim!80!black, font=\footnotesize, below] at (2.55,-0.34) {$\E\{Y\}$};
    \node[figTerm, font=\footnotesize] at (5.0,0.5) {area $=$ probability};
  \end{scope}
\end{tikzpicture}
\end{Fig}

\begin{ObjL}{Concept}{1.4.3}{The problem the next theorem solves}
You know the distribution of $Y$, and you need the expectation not of $Y$ itself but of some
transformed quantity $g(Y)$: $Y^2$ on the way to a variance, a loss $L(\delta(Y),\theta)$ on
the way to a risk, an estimator $(1+X)/(2+n)$ built from a count. Now $g(Y)$ is a random
variable in its own right, and \xref{1.4.1} is unambiguous about what its expectation means:
weight each value $z$ that $g(Y)$ can take by the probability that it takes it---compute
against the distribution \emph{of $g(Y)$}. Taken literally, that instruction opens a
preliminary project: derive that distribution. List the values $g$ can produce; for each one,
find all the outcomes $y$ that $g$ sends there and merge their probabilities; only then
average. The project is real work, and it would have to be redone for every new $g$. The next
theorem says the project is unnecessary: the same number can be read directly off the
distribution you already hold, by weighting $g(y)$, rather than $y$, by $p(y)$. The name is
the profession's joke at its own expense---the formula is what everyone writes down
instinctively, without noticing that, by the definition, there was something to prove; the
theorem is why the instinct is correct. And it earns its keep at once: the risk function of
\xref{1.5.8} \emph{is} this formula with $g$ the loss, a variance is this formula with
$g(y)=(y-\mu)^2$, and its bivariate form,
$\E\{g(X,Y)\}=\sum_x\sum_y g(x,y)\,p(x,y)$, is behind every covariance.
\end{ObjL}

\begin{Obj}{Theorem}{1.4.4}{The law of the unconscious statistician}
Let $Y$ be discrete with pmf $p$, and let $g$ be a real function. Then
\[ \E\{g(Y)\} \;=\; \sum_{y} g(y)\,p(y) \]
(and $\E\{g(Y)\}=\int g(y)f(y)\,dy$ in the continuous case): the expectation of $g(Y)$ may be
computed against the distribution of $Y$, without first finding the distribution of $g(Y)$.
Equivalently: grouping the average by outcome (each $y$ weighted by $p(y)$) and grouping it by
value (each $z$ weighted by $P(g(Y)=z)$) must agree. The theorem is a regrouping identity, and
its proof is exactly that regrouping.
\end{Obj}
\begin{Prf}{1.4.4}{Proof by Regrouping of the Unconscious-Statistician Formula (discrete case).}{}
\item Let $Z=g(Y)$, a random variable with pmf $p_Z(z)=P(g(Y)=z)=\sum_{y:\,g(y)=z}p(y)$.
\item By \xref{1.4.1},
      \[ \E\{Z\} \;=\; \sum_z z\,p_Z(z) \;=\; \sum_z\, z \!\sum_{y:\,g(y)=z}\! p(y) . \]
\item Each outcome $y$ appears in exactly one inner sum, contributing $g(y)\,p(y)$; regrouping
      the double sum by $y$ (legitimate by absolute convergence) gives
      $\E\{g(Y)\}=\sum_y g(y)\,p(y)$.
\end{Prf}

\begin{ObjL}{Example}{1.4.5}{LOTUS on a many-to-one map}
Let $Y$ be uniform on $\{-1,0,1\}$ and $g(y)=y^2$. The definition-first route derives the
distribution of $Z=g(Y)$: the two outcomes $-1$ and $1$ both land on the value $1$, so their
probabilities merge, $P(Z=1)=\tfrac23$, while $P(Z=0)=\tfrac13$; then
$\E\{Z\}=0\cdot\tfrac13+1\cdot\tfrac23=\tfrac23$. The theorem's route never mentions $Z$'s
distribution:
\[ \E\{g(Y)\} \;=\; (-1)^2\cdot\tfrac13+0^2\cdot\tfrac13+1^2\cdot\tfrac13 \;=\; \tfrac23 . \]
Same number, different bookkeeping---and the difference is the whole content of the theorem:
the first route's single term $1\cdot\tfrac23$ is the second route's two terms
$(-1)^2\cdot\tfrac13+1^2\cdot\tfrac13$, regrouped. That merging of outcomes that share a value
is all the proof of \xref{1.4.4} does in general.
\end{ObjL}

\begin{Obj}{Definition}{1.4.6}{The joint distribution of a pair}
For a pair of random variables $(X,Y)$, the joint distribution of \xref{1.1.6} specializes to
a joint pmf $p(x,y)=P(X=x,\,Y=y)$ in the discrete case, or a joint density $f(x,y)$ in the
continuous case.
\end{Obj}

\begin{ObjL}{Remark}{1.4.7}{Reading the joint: the two-way table}
The joint is the complete description of the pair: it carries not only how each variable
behaves alone but every trace of the dependence between them. For discrete variables it is
usefully pictured as a two-way table---one row per value $x$, one column per value $y$, the
cell holding $p(x,y)$---with all the cells summing to $1$. The next two definitions are the
table's row totals and its rescaled rows.
\end{ObjL}

\begin{Obj}{Definition}{1.4.8}{Marginal distribution}
The \emph{marginal} pmf of $X$ recovers the distribution of $X$ alone by summing the joint
over everything that can happen to $Y$:
\[ p_X(x) \;=\; \sum_y p(x,y) \]
(and $f_X(x)=\int f(x,y)\,dy$ in the continuous case).
\end{Obj}

\begin{ObjL}{Remark}{1.4.9}{Why ``marginal,'' and what marginalizing forgets}
In the two-way table the marginal's values are the row totals, classically written in the
table's \emph{margin}---hence the name. And marginalizing is a one-way street: the joint
determines both marginals, but the marginals do not determine the joint, because everything
about the dependence has been summed away (\xref{1.1.6}).
\end{ObjL}

\begin{Obj}{Definition}{1.4.10}{Conditional distribution}
For any $x$ with $p_X(x)>0$, the \emph{conditional} pmf of $Y$ given $X=x$ is
\[ p_{Y\mid X}(y\mid x) \;=\; \frac{p(x,y)}{p_X(x)} \]
(and $f_{Y\mid X}(y\mid x)=f(x,y)/f_X(x)$ for densities, when $f_X(x)>0$).
\end{Obj}

\begin{ObjL}{Remark}{1.4.11}{Conditioning as reweighting a slice}
The recipe behind the formula: fix $X=x$, take that row of the table---the $x$-slice of the
joint---and rescale it by its total so it sums to $1$. The result is the distribution of $Y$
for someone who has learned that $X=x$: conditioning is how a distribution absorbs
information.
\end{ObjL}

\begin{ObjL}{Remark}{1.4.12}{Conditional distributions are probability distributions!}
For each fixed $x$, the numbers $p_{Y\mid X}(y\mid x)$ are nonnegative and sum to $1$: a
conditional distribution is a genuine probability distribution, so everything defined for
distributions---expectations included---applies to it verbatim. That single observation powers
the next two objects.
\end{ObjL}

\begin{Obj}{Definition}{1.4.13}{Conditional expectation}
The \emph{conditional expectation} of $Y$ given $X=x$ is the expectation of the conditional
distribution (\xref{1.4.10}),
\[ \E\{Y\mid X=x\} \;=\; \sum_y y\,p_{Y\mid X}(y\mid x), \]
and $\E\{Y\mid X\}$ denotes the random variable obtained by evaluating this function of $x$ at
$X$.
\end{Obj}

\begin{Obj}{Theorem}{1.4.14}{Iterated expectations}
Whenever the expectations exist,
\[ \E\{\E\{Y\mid X\}\} \;=\; \E\{Y\}. \]
\end{Obj}
\begin{Prf}{1.4.14.a}{Proof by Conditioning of the Iterated-Expectations Identity (discrete case).}{}
\item $\E\{Y\mid X\}$ is the function $x\mapsto\E\{Y\mid X=x\}$ evaluated at $X$, so by the law
      of the unconscious statistician (\xref{1.4.4}),
      \[ \E\{\E\{Y\mid X\}\} \;=\; \sum_x \E\{Y\mid X=x\}\,p_X(x) . \]
\item Substitute the conditional expectation (\xref{1.4.13}),
      $\E\{Y\mid X=x\}=\sum_y y\,p_{Y\mid X}(y\mid x)$:
      \[ \E\{\E\{Y\mid X\}\}
         \;=\; \sum_x\Big(\sum_y y\,p_{Y\mid X}(y\mid x)\Big)\,p_X(x) . \]
\item Each outer term carries the factor $p_X(x)$, so the terms with $p_X(x)=0$ vanish; in the
      terms with $p_X(x)>0$, substitute the conditional pmf (\xref{1.4.10}),
      $p_{Y\mid X}(y\mid x)=p(x,y)/p_X(x)$:
      \[ \E\{\E\{Y\mid X\}\}
         \;=\; \sum_x\Big(\sum_y y\,\frac{p(x,y)}{p_X(x)}\Big)\,p_X(x) . \]
\item The factor $p_X(x)$ does not involve $y$, so it passes inside the inner sum, where it
      cancels the denominator term by term,
      \[ y\,\frac{p(x,y)}{p_X(x)}\cdot p_X(x) \;=\; y\,p(x,y) ,
         \qquad\text{leaving}\qquad
         \E\{\E\{Y\mid X\}\} \;=\; \sum_x\sum_y y\,p(x,y) . \]
\item Exchange the order of summation (legitimate by the assumed absolute convergence), and
      pull $y$---constant in $x$---out of the inner sum:
      \[ \sum_x\sum_y y\,p(x,y) \;=\; \sum_y\sum_x y\,p(x,y) \;=\; \sum_y y\sum_x p(x,y) . \]
\item The inner sum is the marginal pmf of $Y$ (\xref{1.4.8}, with the roles of $X$ and $Y$
      exchanged): $\sum_x p(x,y)=p_Y(y)$. Hence, by the definition of expectation
      (\xref{1.4.1}),
      \[ \E\{\E\{Y\mid X\}\} \;=\; \sum_y y\,p_Y(y) \;=\; \E\{Y\} . \]
\end{Prf}
\begin{Prf}{1.4.14.b}{Proof by Conditioning of the Iterated-Expectations Identity (continuous case).}{}
\item As before, $\E\{Y\mid X\}$ is a function of $X$, so by the law of the unconscious
      statistician (\xref{1.4.4}), now in its integral form,
      \[ \E\{\E\{Y\mid X\}\} \;=\; \int \E\{Y\mid X=x\}\,f_X(x)\,dx . \]
\item Substitute the conditional expectation (\xref{1.4.13}),
      $\E\{Y\mid X=x\}=\int y\,f_{Y\mid X}(y\mid x)\,dy$:
      \[ \E\{\E\{Y\mid X\}\}
         \;=\; \int\!\left(\int y\,f_{Y\mid X}(y\mid x)\,dy\right) f_X(x)\,dx . \]
\item The outer integrand carries the factor $f_X(x)$, so the set where $f_X(x)=0$
      contributes zero to the outer integral; on the set where $f_X(x)>0$, substitute the
      conditional density (\xref{1.4.10}), $f_{Y\mid X}(y\mid x)=f(x,y)/f_X(x)$:
      \[ \E\{\E\{Y\mid X\}\}
         \;=\; \int\!\left(\int y\,\frac{f(x,y)}{f_X(x)}\,dy\right) f_X(x)\,dx . \]
\item The factor $f_X(x)$ does not involve $y$, so it passes inside the inner integral, where
      it cancels the denominator pointwise,
      \[ y\,\frac{f(x,y)}{f_X(x)}\cdot f_X(x) \;=\; y\,f(x,y) ,
         \qquad\text{leaving}\qquad
         \E\{\E\{Y\mid X\}\} \;=\; \int\!\!\int y\,f(x,y)\,dy\,dx . \]
\item Exchange the order of integration (legitimate, by Fubini's theorem, under the assumed
      absolute convergence), and pull $y$---constant in $x$---out of the inner integral:
      \[ \int\!\!\int y\,f(x,y)\,dy\,dx \;=\; \int\!\!\int y\,f(x,y)\,dx\,dy
         \;=\; \int y\!\left(\int f(x,y)\,dx\right) dy . \]
\item The inner integral is the marginal density of $Y$ (\xref{1.4.8}, with the roles of $X$
      and $Y$ exchanged): $\int f(x,y)\,dx=f_Y(y)$. Hence, by the definition of expectation
      (\xref{1.4.1}),
      \[ \E\{\E\{Y\mid X\}\} \;=\; \int y\,f_Y(y)\,dy \;=\; \E\{Y\} . \]
\end{Prf}

% ============================================================
\sub{1.5}{Decision Theory: Actions, Loss, and Risk}

Section \xref{1.3} left the clinical trial in a precise but incomplete state: a model (two
independent Bernoulli samples, $\theta=(\theta_T,\theta_P)\in[0,1]^2$) and an analytic goal
(learn the difference $\theta_T-\theta_P$). What it did not supply is a \emph{method}---and
candidate methods are cheap. Report the difference of the observed rates; report the shaded
difference of \xref{1.8.3}; ignore the data and always report $0$; average the first two. Every
one of them produces an answer on every data set, so producing answers cannot be what
``working'' means. The deck's question---``But first, what do we mean `good?'\,''---is not
rhetorical: until \emph{good} is defined, \emph{best} is meaningless.

Decision theory is the discipline of making ``good'' explicit, and its move is to look past
the answer to its \emph{consequences}. Choosing is acting: report a number, declare efficacy,
approve or shelve. Actions have consequences that depend on the unknown state of
reality---report a difference of $0.30$ when the truth is $0.05$ and patients are steered to a
nearly useless treatment; declare ``no effect'' when there is one and a working therapy is
shelved. So attach to every (action, truth) pair a price: the \emph{loss}. Judging a
\emph{method}---a rule, fixed before the data arrive, committing you to an action for every
data set you might see---then means asking what it can be \emph{expected} to lose. Because the
data are random, that expected loss is computable from the model; because the truth is
unknown, it must be computed at every candidate $\theta$. The result is not a number but a
curve over the parameter space: the \emph{risk}.

A curve is an honest but awkward report card. Risk curves cross---the rule ``always report
$0$'' is unbeatable when $\theta_T=\theta_P$ and dreadful elsewhere---so ``minimize the risk''
cannot be a theorem-backed instruction (a principle, the deck insists, not a result). What
survives is a discipline: discard any method whose curve another beats everywhere
(admissibility), and, when a single number is wanted, average the curve against a prior (the
Bayes risk of \S\,\xref{1.6}). This framing of statistics---estimation and testing as one
decision problem with different action spaces and losses (\S\,\xref{1.8})---is Abraham Wald's,
worked out in the 1940s, and the course adopts it from the outset. The definitions below fix
the vocabulary; each is already visible in the trial. Decision theory begins from the
statistical model of \xref{1.3.1}---a ``parameterized set of distributions on the sample
space''---and adds three objects. \HPS{1}

\begin{Obj}{Definition}{1.5.1}{Action space}
The \emph{action space} $\mathcal{A}$ is the set of possible choices of the decision maker:
the values an analysis is allowed to output.
\end{Obj}

\begin{ObjL}{Remark}{1.5.2}{Choosing the action space}
The action space is fixed by the \emph{question}, not by the data. The same trial supports
estimation---report a plausible value of the cure-rate difference, $\mathcal{A}=[-1,1]$, every
value a difference of two proportions could take---or testing---assert ``efficacious'' or
``not,'' an action space with exactly two elements; changing $\mathcal{A}$ changes the problem
(\S\,\xref{1.8} makes both versions precise). And note what $\mathcal{A}$ is not: it is not
the sample space---data live in $S$, conclusions live in $\mathcal{A}$, and the method of
\xref{1.5.5} will be the bridge---and its elements carry no probabilities, because actions are
\emph{chosen}, not drawn.
\end{ObjL}

\begin{Obj}{Definition}{1.5.3}{Loss function}
The \emph{loss function} is a map
\[ L : \mathcal{A}\times\Theta \;\to\; \mathbb{R}, \]
where $L(a,\theta)$ tells you how much you lose by choosing action $a$ when the parameter is
$\theta$.\note{The deck words the domain both ways---``the parameter space cross the action
space'' on the definition slides, ``the action space cross the parameter space'' in the
typeset cheat sheet---while always writing the arguments $L(a,\theta)$. These notes follow the
cheat sheet.}
\end{Obj}

\begin{ObjL}{Remark}{1.5.4}{Reading the loss function}
The loss is a \emph{modeling choice}, as much a part of the problem's specification as the
statistical model itself: it declares, in advance and in numbers, how bad each mistake would
be---zero (or least) for the best action available at $\theta$, growing as the action worsens,
in whatever units of regret the analyst adopts. The course will stick to two standards
(\S\,\xref{1.8}): \emph{squared-error loss}, $L(a,\theta)=(a-\theta)^2$, for
estimation---near-misses almost free, large misses punished disproportionately (double the
miss, quadruple the loss)---and \emph{zero--one loss} for testing, which charges $1$ for a
wrong assertion and $0$ for a right one. Nothing in the definition demands symmetry---in the
trial, shelving a working therapy may well cost more than approving a useless one, and a loss
function can say so---but the two standards keep the theory clean. For the trial, with the
difference as the target: $L(a,\theta)=\big(a-(\theta_T-\theta_P)\big)^2$.
\end{ObjL}

\begin{Obj}{Definition}{1.5.5}{Method (decision rule)}
A \emph{method}, or \emph{decision rule}, is a function
\[ \delta : S \;\to\; \mathcal{A} \]
from the sample space to the action space: what you will do, specified in advance, for every
data set you might observe.
\end{Obj}

\begin{ObjL}{Remark}{1.5.6}{A plan, not an answer}
The force of the definition is in its timing. A method is chosen \emph{before} the data
arrive, and it commits you on every data set you might see---a complete contingency plan, not
an answer. That is exactly what makes methods judgeable: a single answer cannot be evaluated
until the truth is known, but a plan can be averaged over everything the model says the data
may do. And since the data are modeled as random, $\delta(Y)$ is a random variable---so it has
an expected loss, and that is where the risk comes from.
\end{ObjL}

\begin{ObjL}{Example}{1.5.7}{Three methods for the trial}
Each of the following is a method from $S=\{0,1\}^{n_T}\times\{0,1\}^{n_P}$ to
$\mathcal{A}=[-1,1]$:
\[ \delta_1(y)=\frac{x_T}{n_T}-\frac{x_P}{n_P}, \qquad
   \delta_2(y)=\frac{1+x_T}{2+n_T}-\frac{1+x_P}{2+n_P}, \qquad
   \delta_3(y)\equiv 0 , \]
the difference of observed rates, the shaded version of \xref{1.8.3}, and a rule that ignores
the data entirely. Nothing in the definition requires a method to be sensible---that is the
point: ``method'' is the arena, and the rest of the section builds the referee.
\end{ObjL}

\begin{Obj}{Definition}{1.5.8}{Risk function}
The \emph{risk function} of a decision rule $\delta$ is the expected loss as a function of the
parameter:
\[ R_{\delta}(\theta) \;=\; \E^{\theta}\{L(\delta(Y),\theta)\}
   \;=\; \int_{y\in S} f^{\theta}(y)\,L(\delta(y),\theta)\,dy
   \quad\text{or}\quad \sum_{y\in S} p^{\theta}(y)\,L(\delta(y),\theta), \]
where $\E^{\theta}$ denotes expectation when the data's distribution is the one indexed by
$\theta$.
\end{Obj}

\begin{ObjL}{Remark}{1.5.9}{Reading the risk function}
Read the construction from the inside out: fix a candidate truth $\theta$; the model hands the
data their distribution; the method turns the random data into a random action $\delta(Y)$;
the loss turns the random action into a random number $L(\delta(Y),\theta)$; and the
expectation---the display in \xref{1.5.8} is LOTUS (\xref{1.4.4}) with
$g(y)=L(\delta(y),\theta)$---collapses the randomness to a single number: the method's average
penalty \emph{if $\theta$ were the truth}. Sweeping $\theta$ across $\Theta$ traces the risk
curve (Figure~\xref{1.5.10} shows the machine; Figure~\xref{1.5.13} shows the curves). Two
instances carry the course: under zero--one loss the risk \emph{is} the probability of a wrong
assertion, $R_{\delta}(\theta)=P^{\theta}(\delta\ \text{errs})$; under squared-error loss it
is the mean squared error, which splits as variance plus squared bias (\xref{1.A.1}). The
trial's two estimators have their risks computed in closed form in \xref{1.A.2} and
\xref{1.A.3}.
\end{ObjL}

\begin{Fig}{1.5.10}{The anatomy of a decision problem}{How the four objects compose, at one
fixed candidate truth $\theta$: the model hands the data their distribution, the method maps
data to an action, the loss prices that action against $\theta$, and the expectation over the
data collapses the randomness to the single number $R_{\delta}(\theta)$. Sweeping $\theta$
over $\Theta$ traces the risk curves of Figure~\xref{1.5.13}
(\xref{1.5.1}, \xref{1.5.3}, \xref{1.5.5}, \xref{1.5.8}).}
\begin{tikzpicture}[x=1cm,y=1cm,
  every node/.style={font=\small},
  box/.style={draw, rounded corners=2pt, align=center, inner sep=5pt, minimum height=0.78cm,
    execute at begin node={\hyphenpenalty=10000\exhyphenpenalty=10000}},
  arr/.style={->, thick}]
  \node[box, draw=figLim, fill=figLim!10] (truth) at (0,1.7)
    {\textbf{truth} $\theta\in\Theta$};
  \node[box, draw=cuBlueDk, fill=cuBlueLt!45] (data) at (4.6,1.7)
    {\textbf{data} $Y\in S$};
  \node[box, draw=figTerm, fill=figTerm!8] (act) at (9.2,1.7)
    {\textbf{action} $\delta(Y)\in\mathcal{A}$};
  \node[box, draw=figBad, fill=figBad!8] (loss) at (9.2,0)
    {\textbf{loss} $L(\delta(Y),\theta)$};
  \node[box, draw=figLim, fill=figLim!10] (risk) at (4.6,0)
    {\textbf{risk} $R_{\delta}(\theta)$};
  \draw[arr, figLim!80!black] (truth) -- node[above, font=\footnotesize]{the model $p^{\theta}$} (data);
  \draw[arr, figTerm] (data) -- node[above, font=\footnotesize]{the method $\delta$} (act);
  \draw[arr, figTerm] (act) -- node[right, font=\footnotesize]{priced by $L$} (loss);
  \draw[arr, figLim!80!black] (truth.south) -- (0,-1.15)
    -- node[below, font=\footnotesize]{$\theta$ enters the ledger too} (9.2,-1.15)
    -- (loss.south);
  \draw[arr, figBad!85!black] (loss) -- node[above, font=\footnotesize]{$\E^{\theta}$} (risk);
\end{tikzpicture}
\end{Fig}

\begin{Obj}{Principle}{1.5.11}{The risk principle}
``THIS IS A PRINCIPLE, NOT A RESULT: try to minimize the risk.'' And---``Remember, Risk is a
function!'' One method's risk is an entire curve over $\Theta$, not a number, so two methods
need not be comparable: each may be better somewhere. The principle disciplines the search for
methods; it does not by itself select one.
\end{Obj}

\begin{Obj}{Definition}{1.5.12}{Admissibility}
A method is \emph{admissible} if there is no other method that does at least as well for every
parameter value \emph{and} better for some parameter value. A method that is so dominated is
\emph{inadmissible}: whatever $\theta$ turns out to be, you would not regret abandoning it.
\end{Obj}

\begin{Fig}{1.5.13}{Three risk functions}{Risk curves for three decision rules, as sketched in
lecture. The dashed rule $\delta_3$ lies strictly above $\delta_1$ everywhere, so $\delta_1$
dominates it and $\delta_3$ is inadmissible; $\delta_1$ and $\delta_2$ cross, so neither
dominates the other, and the deck's legend marks both admissible (\xref{1.5.12}).}
\begin{tikzpicture}[x=1.15cm,y=1.9cm]
  % axes
  \draw[->] (-3.3,0) -- (4.55,0) node[right=2pt] {$\theta$};
  \draw[->] (-3,-0.08) -- (-3,1.95) node[left=2pt] {$R(\theta)$};
  \foreach \x in {-2,0,2,4} \draw (\x,0) -- (\x,-0.04) node[below,font=\footnotesize] {$\x$};
  \foreach \y in {0.5,1.0} \draw (-3,\y) -- (-3.06,\y) node[left,font=\footnotesize] {\y};
  % delta_1: blue solid
  \draw[figTerm, thick] plot [smooth, tension=0.9]
    coordinates {(-3,0.80) (-1.5,1.00) (0,0.80) (2,0.55) (3.9,0.75)};
  \node[figTerm, font=\footnotesize, right] at (3.95,0.68) {$\delta_1$};
  % delta_3: red dashed, strictly above delta_1
  \draw[figBad, thick, dashed] plot [smooth, tension=0.9]
    coordinates {(-3,0.89) (-1.5,1.09) (0,0.89) (2,0.64) (3.9,0.84)};
  \node[figBad, font=\footnotesize, right] at (3.95,0.86) {$\delta_3$};
  % delta_2: green solid, crosses the others
  \draw[figLim, thick] plot [smooth, tension=0.9]
    coordinates {(-3,0.50) (-1,0.85) (1.5,1.15) (3.9,0.90)};
  \node[figLim, font=\footnotesize, right] at (3.95,1.06) {$\delta_2$};
  % domination annotation (label above, leader into the delta_1-delta_3 gap, as in the deck)
  \draw[black!70, densely dotted] (-0.5,1.42) -- (-0.5,0.98);
  \node[black!70, font=\footnotesize, above] at (-0.5,1.42)
    {$\delta_1$ dominates $\delta_3$ here};
  % legend (kept clear of the curves, upper right)
  \node[draw=black!40, fill=white, rounded corners=1.5pt, anchor=north east, align=left,
        font=\footnotesize, inner sep=4pt] at (4.25,1.95)
    {\textcolor{figTerm}{\rule[2pt]{14pt}{1.2pt}}\; $\delta_1$ --- admissible\\
     \textcolor{figLim}{\rule[2pt]{14pt}{1.2pt}}\; $\delta_2$ --- admissible\\
     \textcolor{figBad}{\rule[2pt]{4pt}{1.2pt}\,\rule[2pt]{4pt}{1.2pt}\,\rule[2pt]{4pt}{1.2pt}}\;
       $\delta_3$ --- inadmissible};
\end{tikzpicture}
\end{Fig}

\begin{ObjL}{Concept}{1.5.14}{Should we follow the normative prescription?}
If you have a coherent ordering on decision rules---any consistent way of ranking them---then
that ordering is consistent with the risk for \emph{some} loss function. ``And so why not use
THE loss function?!'' Ranking methods at all commits you, implicitly, to behaving as if you
were minimizing some risk; decision theory just makes the loss explicit.\note{Stated at the
deck's informal level. A precise version---axioms of coherence and the accompanying
complete-class theorems---is beyond this course.}
\end{ObjL}

% ============================================================
\sub{1.6}{Bayesian Decision Theory}

Bayesian decision theory keeps the statistical model, the action space, and the loss function
of \S\,\xref{1.5}, and adds one object. \HPS{1}

\begin{Obj}{Definition}{1.6.1}{Prior distribution}
The \emph{prior distribution} $\pi$ is a probability distribution on the parameter space
$\Theta$: the parameter is treated as random, with $\pi$ describing where it is likely to be
before the data are seen.
\end{Obj}

\begin{Obj}{Definition}{1.6.2}{Bayes risk}
The \emph{Bayes risk} of a decision rule $\delta$ under the prior $\pi$ is the expected
risk---``treating the parameter as random'':
\[ \E\{R_{\delta}(\theta)\} \;=\; \int_{\theta\in\Theta} \pi(\theta)\,R_{\delta}(\theta)\,d\theta \]
(a sum, for a discrete prior).\note{The deck prints this integral without the trailing
$d\theta$; it is restored here.} Where the risk was a function on $\Theta$, the Bayes risk is a
single number: the prior averages the curve.
\end{Obj}

\begin{Obj}{Principle}{1.6.3}{The Bayes principle}
The Bayes version of the principle: minimize the \emph{expected} risk. Among all decision
rules, choose one whose Bayes risk is smallest---a \emph{Bayes rule} for the prior $\pi$.
\end{Obj}

\begin{ObjL}{Concept}{1.6.4}{Where does the prior distribution come from?}
Maybe the parameter really is random and you know its distribution---or at least have some data
to make a good guess about it. Or maybe you just make the prior up: ``it can be helpful!''
\end{ObjL}

\begin{ObjL}{Remark}{1.6.5}{Should we follow the Bayes prescription?}
If we believe in the prior, then there really is only one distribution for the data, and if we
believe the general prescription (minimize risk), we should minimize the corresponding
risk---the Bayes risk. Also: ``(almost) every admissible rule is a Bayes rule and every Bayes
rule is admissible\dots\ so even if we don't have a prior, the Bayes rule is
sensible.''\note{As stated in lecture, with the instructor's ``(almost).'' The two directions
require side conditions in general; the precise statements are the complete-class theorems,
beyond this course.}
\end{ObjL}

% ============================================================
\sub{1.7}{The Cheat Sheet: Decision Theory in Eleven Items}

\begin{ObjL}{Concept}{1.7.1}{The essential points}
The instructor's own typeset summary of \S\S\,\xref{1.3}--\xref{1.6}, reproduced with the
notation of these notes:
\begin{enumerate}[leftmargin=2.2em,itemsep=3pt,topsep=3pt]
  \item Sample space $S$.
  \item Action space $\mathcal{A}$: the possible actions available to the decision-maker.
  \item Parameter space $\Theta$.
  \item Data $Y$ taking values in $S$.
  \item There is a statistical model: for each $\theta$, a distribution for $Y$, with
        $f^{\theta}$ or $p^{\theta}$ the corresponding density or pmf.
  \item Loss function $L:\mathcal{A}\times\Theta\to\mathbb{R}$: $L(a,\theta)$ is the loss
        associated with making decision $a$ when the true distribution is the one with pmf or
        density $p^{\theta}$ or $f^{\theta}$.
  \item A decision rule $\delta$ maps the sample space to the action space.
  \item The risk function of $\delta$, $R_{\delta}(\theta)$, maps each $\theta\in\Theta$ to the
        expected loss when the decision rule is $\delta$ and the distribution of $Y$ is the one
        corresponding to $\theta$:
        \[ R_{\delta}(\theta)=\E^{\theta}\{L(\delta(Y),\theta)\}
           =\int_{y\in S} f^{\theta}(y)\,L(\delta(y),\theta)\,dy
           \ \ \text{or}\ \ \sum_{y\in S} p^{\theta}(y)\,L(\delta(y),\theta). \]
  \item Normative principle of decision theory: use only admissible decision rules---there must
        be no other rule at least as good for every $\theta$ and better for some $\theta$.
  \item Prior distribution $\pi$ for a random $\theta$ taking values in $\Theta$; the Bayes
        risk of $\delta$ under $\pi$ is
        \[ \E\{R_{\delta}(\theta)\}=\int_{\theta\in\Theta}\pi(\theta)\,R_{\delta}(\theta)\,d\theta . \]
  \item The Bayesian principle: choose $\delta$ to minimize the Bayes risk.
\end{enumerate}
\end{ObjL}

% ============================================================
\sub{1.8}{Estimation, Testing, and the Clinical Trial}

Two special cases of the framework carry the whole course. \HPS{1}

\begin{Obj}{Definition}{1.8.1}{Estimation}
\emph{Estimation} is the special case in which the action space is the parameter space (or the
set of values of the quantity being estimated\note{The deck says exactly ``the action space is
the parameter space''; the parenthetical broadening is ours, and is needed as soon as one
estimates a \emph{function} of the parameter---in the trial of \xref{1.8.4},
$\mathcal{A}=[-1,1]$ while $\Theta=[0,1]^2$.}): the decision \emph{is} a guess at the
parameter. The loss function reflects a distance between the parameter and the estimate; this
course will stick to \emph{squared-error loss},
\[ L(a,\theta) \;=\; (a-\theta)^2 . \]
\end{Obj}

\begin{Obj}{Definition}{1.8.2}{Testing}
\emph{Testing} is the special case in which the action space has two choices: assert that the
parameter is in one portion of the parameter space, or that it is in the rest. The loss
function reflects whether you guessed right; this course will stick to \emph{zero--one loss},
\[ L(a,\theta) \;=\; \one\{a \text{ is the wrong assertion about } \theta\} , \]
which loses $1$ for a wrong call and $0$ for a right one.
\end{Obj}

\begin{ObjL}{Example}{1.8.3}{Two Bernoulli samples, and two estimators}
The running example: two samples of independent Bernoulli outcomes with common expectations
within each sample---$\theta_T$ in one arm, $\theta_P$ in the other. Estimate the difference
$\theta_T-\theta_P$ between the two expectations, with squared-error loss. The Bayes version
puts a Beta$(1,1)$ (uniform) prior on each parameter, the two assumed independent. Two
estimators will be compared:
\begin{enumerate}[leftmargin=2.2em,itemsep=2pt,topsep=2pt]
  \item the difference between the two observed rates, $X_T/n_T - X_P/n_P$; and
  \item the difference if, separately in each sample, one instead uses
        \[ \frac{1+X}{2+n}, \]
        where $X$ is the sum of the Bernoulli variables and $n$ is the sample size.
\end{enumerate}
Here the second estimator shades the observed rate toward $\tfrac12$; where it comes from, and
when it wins, is the business of the appendix (\appref{1.A.3}, \appref{1.A.4}).
\end{ObjL}

\begin{ObjL}{Example}{1.8.4}{The clinical trial, formalized}
For the two-arm trial of \xref{1.3.4}, with $n_T$ treatment patients and $n_P$ placebo
patients and cure ($1$) or failure ($0$) recorded for each:
\begin{itemize}[leftmargin=1.8em,itemsep=2pt,topsep=2pt]
  \item \emph{Sample space:} $S=\{0,1\}^{n_T}\times\{0,1\}^{n_P}$, the possible cure/failure
        records.
  \item \emph{Statistical model and parameter space:} within each arm the outcomes are
        independent Bernoulli variables with a common expectation; the parameter is
        $\theta=(\theta_T,\theta_P)$ with $\Theta=[0,1]^2$, and
        \[ p^{\theta}(y)=\theta_T^{\,x_T}(1-\theta_T)^{n_T-x_T}\,
                       \theta_P^{\,x_P}(1-\theta_P)^{n_P-x_P}, \]
        where $x_T,x_P$ are the arm totals of $y$.
  \item \emph{Analytic goal:} learn the difference $\theta_T-\theta_P$ (is the treatment
        efficacious?).
  \item \emph{Action space:} the possible values of the difference, $\mathcal{A}=[-1,1]$.
  \item \emph{Loss function:} squared error, $L(a,\theta)=\big(a-(\theta_T-\theta_P)\big)^2$.
\end{itemize}
The lecture closes by naming, for each of the two estimators of \xref{1.8.3}: the risk
function, the prior, and the Bayes risk. The deck sets these up and leaves the computations to
us; they are worked in the appendix (\appref{1.A.2}--\appref{1.A.4}).
\end{ObjL}

\begin{ObjL}{Remark}{1.8.5}{The probability theory we will need}
The deck lists ``Bias--variance decomposition'' and ``Beta distributions,'' adding: ``Also,
variances, covariance, rules for expectations and variances for linear combinations of random
variables; And the role of independence\dots'' For the Beta$(\alpha,\beta)$ distribution
($\alpha,\beta>0$) the deck displays the standard
reference facts:\note{The deck shows the pdf/mean/variance rows of a standard reference table;
they are typeset here directly.}
\[ f(x)=\frac{x^{\alpha-1}(1-x)^{\beta-1}}{\mathrm{B}(\alpha,\beta)}
   \quad\text{on } [0,1], \qquad
   \E\{X\}=\frac{\alpha}{\alpha+\beta}, \qquad
   \Var(X)=\frac{\alpha\beta}{(\alpha+\beta)^2(\alpha+\beta+1)} , \]
where $\mathrm{B}(\alpha,\beta)=\Gamma(\alpha)\Gamma(\beta)/\Gamma(\alpha+\beta)$. Note that
Beta$(1,1)$ is the uniform distribution on $[0,1]$.
\end{ObjL}

\begin{ObjL}{Remark}{1.8.6}{What about AI?}
The instructor's advice: set up the calculations yourself---``that is practicing the conceptual
issues, and that is the point''---while acknowledging that hand fluidity with such calculations
``is already pretty much obsolete.'' Two questions are left open: do you need to do some by
hand to understand the conceptual side? And do you need to learn, now, how to get AI to do them
for you?
\end{ObjL}

% ============================================================
\lectureappendix

The deck names the bias--variance decomposition, the risk functions of the two estimators, and
their Bayes risks (Slides 40--41 and the closing checklist), but works none of them. They are
worked here, beyond the deck, in the deck's own notation. Throughout, $X\sim$
Binomial$(n,\theta)$ denotes an arm total, so $\E\{X\}=n\theta$ and
$\Var(X)=n\theta(1-\theta)$.

\begin{ObjL}{Supplement}{1.A.1}{The bias--variance decomposition}
Let $\delta=\delta(Y)$ estimate a target $g(\theta)$ under squared-error loss, with
$\E^{\theta}\{\delta^2\}<\infty$, and write
$m(\theta)=\E^{\theta}\{\delta\}$ and $b(\theta)=m(\theta)-g(\theta)$ (the \emph{bias}). Then
\[ R_{\delta}(\theta) \;=\; \E^{\theta}\{(\delta-g(\theta))^2\}
   \;=\; \Var^{\theta}(\delta) \;+\; b(\theta)^2 . \]
\end{ObjL}
\begin{Prf}{1.A.1}{Proof of the Bias--Variance Decomposition by Adding and Subtracting the Mean.}{}
\item Add and subtract $m(\theta)$ inside the square:
      $(\delta-g)^2=(\delta-m)^2+2(\delta-m)(m-g)+(m-g)^2$.
\item Take $\E^{\theta}$ term by term (linearity). The middle term is
      $2(m-g)\,\E^{\theta}\{\delta-m\}=0$, since $m-g$ is a constant and
      $\E^{\theta}\{\delta\}=m$.
\item What remains is $\E^{\theta}\{(\delta-m)^2\}+(m-g)^2=\Var^{\theta}(\delta)+b(\theta)^2$.
\end{Prf}

\begin{ObjL}{Supplement}{1.A.2}{Risk of the difference of observed rates}
For $\hat\delta = X_T/n_T - X_P/n_P$ in the model of \xref{1.8.4}:
\[ R_{\hat\delta}(\theta)
   \;=\; \frac{\theta_T(1-\theta_T)}{n_T} + \frac{\theta_P(1-\theta_P)}{n_P} . \]
\end{ObjL}
\begin{Prf}{1.A.2}{Proof of the Rate-Difference Risk via Unbiasedness and Independence.}{Uses \xref{1.A.1}.}
\item Each arm's rate is unbiased: $\E^{\theta}\{X_T/n_T\}=n_T\theta_T/n_T=\theta_T$
      (linearity), and likewise for the placebo arm; so
      $b(\theta)=\E^{\theta}\{\hat\delta\}-(\theta_T-\theta_P)=0$.
\item A single Bernoulli variable $Y$ has $\E\{Y^2\}=\E\{Y\}=\theta$, so
      $\Var(Y)=\theta-\theta^2=\theta(1-\theta)$; summing $n$ independent copies and scaling,
      $\Var(X/n)=\theta(1-\theta)/n$ for each arm.
\item The arms are independent, so the variances of the difference add:
      $\Var^{\theta}(\hat\delta)=\theta_T(1-\theta_T)/n_T+\theta_P(1-\theta_P)/n_P$.
\item By the bias--variance decomposition (\xref{1.A.1}) with zero bias, the risk is the
      variance.
\end{Prf}

\begin{ObjL}{Supplement}{1.A.3}{Risk of the shaded estimator}
Write $T=\dfrac{1+X}{2+n}$ for one arm with $X\sim$ Binomial$(n,\theta)$. Then
\[ \E^{\theta}\{T\}=\frac{1+n\theta}{n+2},\qquad
   b(\theta)=\frac{1-2\theta}{n+2},\qquad
   \Var^{\theta}(T)=\frac{n\theta(1-\theta)}{(n+2)^2}, \]
so the one-arm risk is
\[ \mathrm{MSE}_n(\theta) \;=\; \frac{n\theta(1-\theta)+(1-2\theta)^2}{(n+2)^2} , \]
and for $\tilde\delta = T_T - T_P$ (the difference of the shaded rates, arms independent),
\[ R_{\tilde\delta}(\theta) \;=\;
   \frac{n_T\theta_T(1-\theta_T)}{(n_T+2)^2}+\frac{n_P\theta_P(1-\theta_P)}{(n_P+2)^2}
   +\left(\frac{1-2\theta_T}{n_T+2}-\frac{1-2\theta_P}{n_P+2}\right)^{\!2}. \]
At $\theta=\tfrac12$ the one-arm risk is $\tfrac{n}{4(n+2)^2}<\tfrac{1}{4n}$: shading toward
$\tfrac12$ beats the raw rate in the middle of the parameter space, at the price of bias near
the edges---neither estimator dominates the other.
\end{ObjL}
\begin{Prf}{1.A.3}{Proof of the Shaded Estimator's Risk by Linear-Combination Rules.}{Uses \xref{1.A.1}, \xref{1.A.2}.}
\item By linearity, $\E^{\theta}\{T\}=(1+\E^{\theta}\{X\})/(n+2)=(1+n\theta)/(n+2)$, so
      \[ b(\theta)=\frac{1+n\theta}{n+2}-\theta=\frac{1+n\theta-(n+2)\theta}{n+2}
      =\frac{1-2\theta}{n+2} . \]
\item Adding constants does not change variance, and scaling squares:
      $\Var^{\theta}(T)=\Var(X)/(n+2)^2=n\theta(1-\theta)/(n+2)^2$ (step (2) of \xref{1.A.2}).
\item The one-arm risk is variance plus squared bias (\xref{1.A.1}).
\item For the difference, the arms are independent, so variances add, while the bias of a
      difference is the difference of the biases; \xref{1.A.1} assembles the displayed risk.
\item At $\theta=\tfrac12$: bias $=0$ and the risk is $\tfrac{n/4}{(n+2)^2}$, smaller than the
      raw rate's $\tfrac{1}{4n}$ since $n^2<(n+2)^2$.
\end{Prf}

\begin{ObjL}{Supplement}{1.A.4}{Bayes risks under the uniform prior, and where $(1+X)/(2+n)$ comes from}
Under the Beta$(1,1)$ (uniform) prior on each arm's parameter, independently:
\begin{enumerate}[leftmargin=2.2em,itemsep=2pt,topsep=2pt]
  \item the one-arm Bayes risk of the observed rate $X/n$ is
        \[ \int_0^1 \frac{\theta(1-\theta)}{n}\,d\theta=\frac{1}{6n} \,; \]
  \item the one-arm Bayes risk of $T=(1+X)/(2+n)$ is $1/\big(6(n+2)\big)$ --- strictly
        smaller;
  \item for the difference, the Bayes risks add across arms:
        \[ \frac{1}{6n_T}+\frac{1}{6n_P} \ \text{ for the rates,}\quad\text{versus}\quad
           \frac{1}{6(n_T+2)}+\frac{1}{6(n_P+2)} \ \text{ for the shaded estimators;} \]
  \item and $T$ is no accident: under the uniform prior, the posterior distribution of
        $\theta$ given $X=x$ is Beta$(1+x,\,1+n-x)$, whose mean is exactly $(1+x)/(n+2)$ ---
        the shaded estimator is the \emph{Bayes rule} for squared-error loss. (And since the
        arms' posteriors are independent, $T_T-T_P$ is the Bayes rule for the difference.)
\end{enumerate}
\end{ObjL}
\begin{Prf}{1.A.4}{Proof of the Bayes Risks by Integration, and of the Bayes Rule by Identifying the Posterior.}{Uses \xref{1.A.2}, \xref{1.A.3}, \xref{1.8.5}.}
\item $\int_0^1\theta(1-\theta)\,d\theta=\tfrac12-\tfrac13=\tfrac16$, so the rate's risk
      $\theta(1-\theta)/n$ integrates to $1/(6n)$.
\item For $T$: $\int_0^1(1-2\theta)^2\,d\theta=\int_0^1(1-4\theta+4\theta^2)\,d\theta
      =1-2+\tfrac43=\tfrac13$, so
      \[ \int_0^1 \mathrm{MSE}_n(\theta)\,d\theta
      =\frac{n\cdot\frac16+\frac13}{(n+2)^2}
      =\frac{n+2}{6(n+2)^2}=\frac{1}{6(n+2)} . \]
\item For the difference, expand the risk of \xref{1.A.3} and integrate under the independent
      uniform priors. The cross term of the squared bias is
      \[ -2\int_0^1\!\!\int_0^1 \frac{(1-2\theta_T)(1-2\theta_P)}{(n_T+2)(n_P+2)}
      \,d\theta_T\,d\theta_P=0, \]
      because $\int_0^1(1-2\theta)\,d\theta=0$; what survives is the
      sum of the two one-arm Bayes risks (for the rates, the biases are identically zero).
\item For the posterior: with prior density $1$ on $[0,1]$ and $P(X=x\mid\theta)
      =\binom{n}{x}\theta^x(1-\theta)^{n-x}$, the posterior density is proportional to
      $\theta^{x}(1-\theta)^{n-x}=\theta^{(1+x)-1}(1-\theta)^{(1+n-x)-1}$ --- the
      Beta$(1+x,1+n-x)$ kernel. Its mean, by the Beta facts of \xref{1.8.5}, is
      $(1+x)/(n+2)$.
\item Finally, the action $a$ minimizing the posterior expected loss
      $\E\{(\theta-a)^2\mid X=x\}=\Var(\theta\mid X=x)+(\E\{\theta\mid X=x\}-a)^2$
      (by \xref{1.A.1} applied to the posterior) is the posterior mean; and since the Bayes
      risk is the expectation over $X$ of the posterior expected loss (iterated expectations,
      \xref{1.4.14}), minimizing the latter for each $x$ minimizes the former, so $T$ is the
      Bayes rule.
\end{Prf}

\ifSubfilesClassLoaded{\clearpage\objindex}{}

\end{document}
